\scanpage{130}
\[
\tilde r_1^2(\unsure{\tilde u})=\tilde u
\]
\[
\tilde r_1^2(\unsure{\tilde v|\tilde w})=\tilde u^{-1}[\unsure{\tilde u|\tilde w}\tilde v]=\tilde v
\]
\[
\tilde r_1\ell
\]
\[
\deleted{\tilde u^{-1}}\deleted{\tilde u}\tilde v\tilde u\tilde v^{-1}\tilde u^{-1}\;-
\]

\[
(\gamma,\ell)(\gamma',\ell')(\gamma'',\ell'')
\]
\diagram{\begin{tikzcd}[column sep=large,row sep=large]
a \arrow[r,"\ell''"] & \gamma''{}^{-1}(a) \arrow[r,"\gamma''{}^{-1}(\ell')"] & (\gamma'\gamma'')^{-1}(a) \arrow[d,"(\gamma'\gamma'')^{-1}(\ell)"']\\
& & (\gamma\gamma'\gamma'')^{-1}(a)
\end{tikzcd}}

UNIVERSITE DES SCIENCES\\
ET TECHNIQUES DU LANGUEDOC\\
M A T H E M A T I Q U E S\hfill 30.10.80

\begin{center}
SEMINAIRE GEOMETRIE ALGEBRIQUE
\end{center}

Tous les mercredis à 14 h 30, Salle 331.

\begin{tabular}{lll}
Exposés prévus : & 12 novembre, & Schémas formels (C. CONTOU-CARRERE)\\
& 19 novembre, & Schémas formels (suite)\\
& 26 novembre, & Correspondances sur les courbes (D. ALIBERT)\\
& 3 décembre, & Correspondances sur les courbes (suite).
\end{tabular}

\begin{flushright}
D. ALIBERT\qquad C. CONTOU-CARRERE
\end{flushright}

(la première séance aura lieu le 12 novembre).
\margin{130}
